\documentclass[../BTL_PTIT.tex]{subfiles}
\begin{document}

\begin{center}
    \Large{\textbf{TÓM TẮT NỘI DUNG DỰ ÁN}}\\
\end{center}
\vspace{1cm}
\hspace*{15pt}Dự án xây dựng hệ thống giám sát và xử lý chất lượng không khí trong nhà máy theo thời gian thực sử dụng vi điều khiển ESP32, cảm biến khí MQ-135, module ADC ADS1115 và nền tảng giám sát từ xa qua MQTT. Bài toán xuất phát từ nhu cầu phát hiện sớm rò rỉ hoặc gia tăng nồng độ khí độc trong môi trường sản xuất để kịp thời cảnh báo và xử lý, giảm rủi ro mất an toàn lao động.

Nội dung chính của dự án tập trung vào các khối chức năng cốt lõi gồm thu thập dữ liệu nồng độ khí từ hai khu vực giám sát, xử lý dữ liệu tại thiết bị biên, truyền thông lên HiveMQ Broker và hiển thị trạng thái vận hành trên dashboard. Hệ thống hỗ trợ hai chế độ điều khiển AUTO và MANUAL, đồng thời bổ sung Emergency Mode bằng 3 nút bấm vật lý để can thiệp khẩn cấp tại hiện trường. Về phần cứng, hệ thống sử dụng thêm module ADC ADS1115 16-bit để tăng độ chính xác khi đọc tín hiệu analog, servo SG90 để điều khiển góc mở cửa sổ thông gió, còi SFM-27 để cảnh báo âm thanh, cùng nguồn dự phòng pin lithium 18650, mạch chuyển nguồn YX850 và module LM2596 nhằm duy trì hoạt động liên tục khi mất điện lưới.

Kết quả thử nghiệm cho thấy hệ thống phản hồi nhanh khi nồng độ khí tăng cao, dữ liệu cập nhật liên tục với độ trễ thấp và khả năng điều khiển từ xa ổn định. Hệ thống đã triển khai cập nhật firmware từ xa (OTA), cơ chế Watchdog Timer và khôi phục trạng thái sau reset để tăng độ tin cậy vận hành. Trên cơ sở đó, báo cáo đánh giá mức độ đáp ứng các yêu cầu chức năng và phi chức năng, đồng thời đề xuất các hướng phát triển như tăng cường bảo mật TLS đầy đủ và bổ sung mô-đun AI tóm tắt xu hướng ô nhiễm.

\end{document}